\documentclass[conference]{IEEEtran}
\IEEEoverridecommandlockouts

% Packages
\usepackage{cite}
\usepackage{amsmath,amssymb,amsfonts}
\usepackage{algorithmic}
\usepackage{graphicx}
\usepackage{textcomp}
\usepackage{xcolor}
\usepackage{booktabs}
\usepackage{tabularx}
\usepackage{multirow}
\usepackage{url}
\usepackage{microtype}
\usepackage[hidelinks]{hyperref}

% Table column definitions for tabularx
\newcolumntype{Y}{>{\raggedright\arraybackslash}X}
\newcolumntype{s}{>{\hsize=.35\hsize\raggedright\arraybackslash}X}
\newcolumntype{m}{>{\hsize=.8\hsize\raggedright\arraybackslash}X}
\newcolumntype{b}{>{\hsize=1.2\hsize\raggedright\arraybackslash}X}

\begin{document}

\title{AI-Powered News Translation from Kannada to English}

\author{
\IEEEauthorblockN{K S Rithesh Hebbar}
\IEEEauthorblockA{\textit{Dept. of Information Science and Engineering} \\
\textit{M. S. Ramaiah Institute of Technology}\\
Bengaluru, Karnataka, India \\
USN: 1MS23IS049}
\and
\IEEEauthorblockN{Divvela Srihari}
\IEEEauthorblockA{\textit{Dept. of Information Science and Engineering} \\
\textit{M. S. Ramaiah Institute of Technology}\\
Bengaluru, Karnataka, India \\
USN: 1MS23IS039}
\and
\IEEEauthorblockN{K Manjunath}
\IEEEauthorblockA{\textit{Dept. of Information Science and Engineering} \\
\textit{M. S. Ramaiah Institute of Technology}\\
Bengaluru, Karnataka, India \\
USN: 1MS23IS048}
}

\maketitle

\begin{abstract}
Language barriers can limit access to news and information published in regional languages. Kannada is widely used in newspapers, digital news portals, and public information sources across Karnataka. However, readers who are unfamiliar with Kannada may experience difficulties understanding news published in the language. This project proposes an AI-powered web application for translating Kannada news text into English using Natural Language Processing (NLP) and a pretrained multilingual neural machine translation model. IndicTrans2, developed by AI4Bharat, is considered as the proposed translation model because it supports translation involving Kannada and English. The application is intended to provide a web-based interface through which users can enter Kannada news content and obtain its English translation. The system aims to preserve the meaning and context of the source text while producing clear and readable English output. Translation quality can be evaluated using reference-based metrics and human assessment of fluency, adequacy, and meaning preservation. As implementation and evaluation are pending, no experimental performance results are claimed. The proposed project aims to improve access to regional news and demonstrate the application of multilingual Artificial Intelligence in reducing language barriers.
\end{abstract}

\begin{IEEEkeywords}
Kannada-to-English Translation, Generative AI, Natural Language Processing, IndicTrans2, Hugging Face, Neural Machine Translation, Web Application.
\end{IEEEkeywords}

% -----------------------------------------------------------------------------
% SECTION I: INTRODUCTION
% -----------------------------------------------------------------------------
\section{Introduction}
\label{sec:introduction}

Access to timely and reliable news reporting is vital for public awareness, civic participation, and regional governance. In multilingual regions such as Karnataka, India, critical socio-political, economic, and cultural developments are documented primarily in Kannada across print daily newspapers, digital publications, official gazettes, and e-papers. However, residents unfamiliar with the Kannada script, interstate migrants, students, and administrative personnel often face severe language barriers in comprehending these local news reports. 

Automating the translation of Kannada news text into English presents substantial computational and linguistic challenges. Kannada is an agglutinative, morphologically rich Dravidian language with a Subject-Object-Verb (SOV) grammatical structure, contrasting sharply with the analytic Subject-Verb-Object (SVO) typography of English. In addition, Kannada journalistic prose incorporates complex morphophonemic compounds (\textit{sandhi} and \textit{samasa}) and domain-specific administrative terms. While recent Large Language Models (LLMs) demonstrate impressive generative fluency, applying unconstrained generative models directly to regional vernacular news frequently leads to factual hallucinations, distorted named entities, and an inability to trace claims back to original source sentences.

To resolve these challenges, this project introduces a verified, end-to-end web system built upon a modular ``Translate-Then-Summarize'' architectural paradigm. By decoupling sentence-level neural machine translation from grounded abstractive summarization, the system guarantees that all translated content maintains strict $1$-to-$1$ sentence alignment with the source text. The application incorporates automated multi-modal news ingestion (covering live RSS feeds, XML sitemaps, digital PDFs, and e-paper image OCR), multi-tier quality estimation guardrails, and interactive accessibility features including sentence-synchronized audio reading and an in-context conversational Q\&A interface.

% -----------------------------------------------------------------------------
% SECTION II: LITERATURE REVIEW
% -----------------------------------------------------------------------------
\section{Literature Review}
\label{sec:literature_review}

This section reviews ten works that shape the design and evaluation of a Kannada-to-English news translation system. They are grouped into five themes: the neural foundations of machine translation (MT) and the low-resource setting, multilingual pretrained models, Indic-language resources and models, benchmarks and evaluation, and the use of large language models (LLMs) for translation.

\subsection{Foundations and the Low-Resource Setting}
Vaswani et al. [1] introduced the Transformer, a sequence-to-sequence architecture built entirely on attention, with no recurrence or convolution. It reached 28.4 BLEU on the WMT 2014 English-German task and trains more efficiently than earlier recurrent systems. Nearly all modern translation models, including those considered in this work, build on this design. Self-attention relates every source token to every other, which is useful for Kannada-to-English translation, where the subject-object-verb order of Kannada must be rearranged into English subject-verb-object order.

Haddow et al. [2] survey low-resource MT. They describe how such language pairs suffer from limited parallel data, out-of-domain data and noisy data, and they review remedies including web crawling and data creation, use of monolingual data (such as back-translation), multilingual transfer and suitable model choices. The survey frames the main challenges for Kannada, which has far less parallel text than English-centric high-resource pairs, and it motivates reusing multilingual pretrained models rather than training from scratch.

\subsection{Multilingual Pretraining and Massively Multilingual MT}
Liu et al. [3] proposed mBART, a sequence-to-sequence model pretrained by denoising large monolingual corpora in many languages and then fine-tuned for translation. The pretraining gave large gains when parallel data was scarce, which established pretrain-then-fine-tune as a standard approach for low-resource MT. Tang et al. [4] extended this idea to mBART-50 by further pretraining it on 50 languages and fine-tuning it for many-to-many translation. These models are widely available and easy to use, but to our knowledge Kannada is not among the 50 languages mBART-50 covers, which limits their direct suitability for this project.

The NLLB Team [5] scaled multilingual MT to 200 languages. Their work combined large-scale mined bitext, a sparsely gated mixture-of-experts model, and a new evaluation benchmark, FLORES-200, with extra attention to low-resource languages and to safety and quality checks. NLLB models support Kannada and are a strong, openly available general-purpose baseline. They are trained for broad coverage rather than for any single region or genre, so news-specific vocabulary and Indian proper nouns may be handled less well than by region-focused systems.

\subsection{Resources and Models for Indic Languages}
Ramesh et al. [6] released Samanantar, then the largest publicly available parallel corpus for Indic languages, with about 49.7 million English sentence pairs across 11 Indic languages. Roughly 12.4 million pairs were collected from existing sources and 37.4 million were newly mined from the web, using document OCR, multilingual sentence representations and approximate nearest-neighbour search. Models trained on the corpus outperformed earlier systems on public benchmarks such as FLORES. The work shows that data quantity and alignment quality, rather than architecture alone, limit translation for languages like Kannada.

Gala et al. [7] present IndicTrans2, the first open model supporting all 22 scheduled Indian languages, including Kannada. The paper also releases the Bharat Parallel Corpus Collection (BPCC), with 230 million bitext pairs, and IN22, an n-way benchmark with diverse domains, conversational content and India-centric text. The authors report that IndicTrans2 outperforms existing models on several benchmarks, and all models and data are released under permissive licenses. IndicTrans2 is the proposed translation model of this project because of its explicit Kannada support, India-focused data and open availability.

\subsection{Benchmarks and Evaluation}
Goyal et al. [8] introduced the FLORES-101 benchmark, which contains professionally translated sentences from Wikipedia-derived sources in 101 languages. Because the same content is translated into every language, results can be compared across languages and systems, and low-resource directions are evaluated on equal terms. Its sentences are drawn from encyclopedic and general text, so it does not fully represent the headlines, quotations and named-entity-heavy style of news, which argues for supplementing it with news-domain test data.

Rei et al. [9] proposed COMET, a neural evaluation framework trained to predict human judgments of translation quality using the source sentence, the machine output and a reference. COMET correlated better with human ratings than overlap-based metrics such as BLEU in the authors' experiments. Learned metrics are therefore useful alongside BLEU and chrF++ when judging output from morphologically rich languages, although their reliability for Kannada should be confirmed against the human assessment of fluency and adequacy planned in this project.

\subsection{Large Language Models for Translation}
Hendy et al. [10] evaluated GPT models for MT across 18 translation directions, covering high- and low-resource languages. They found that GPT models reach very competitive quality for high-resource languages but have limited capability for low-resource ones. They also report that combining GPT models with conventional MT systems can improve quality further. These findings support using a dedicated, Kannada-aware translation model as the core of the system, with LLMs better suited to supporting tasks such as post-editing or summarisation.

\subsection{Summary and Research Gap}
Table I summarises the reviewed works. Taken together, they show that Transformer-based multilingual models, large Indic corpora and open Indic-specific systems have made Kannada-to-English translation practical, and that general benchmarks and learned metrics make quality measurable. However, the reviewed benchmarks are mostly general-domain and none focuses on Kannada news, so questions about named entities, headlines and news terminology remain open. This project addresses that gap by building a practical web application around IndicTrans2 and assessing it on Kannada news text using automatic metrics and human evaluation.

% -----------------------------------------------------------------------------
% TABLE I: EXTENDED COMPARATIVE LITERATURE REVIEW MATRIX
% -----------------------------------------------------------------------------
\begin{table*}[t]
\caption{Comprehensive Comparison of Reviewed Literature: Methodologies, Scope, Contributions, Drawbacks, and Addressed Gaps}
\label{tab:literature_review}
\centering
\small
\begin{tabularx}{\textwidth}{@{} c s m m b b @{}}
\toprule
\textbf{Ref.} & \textbf{Author(s) \& Year} & \textbf{Architecture / Technique} & \textbf{Language Scope \& Datasets} & \textbf{Core Strengths \& Contributions} & \textbf{Limitations, Drawbacks \& Gaps Addressed} \\
\midrule

\cite{vaswani2017attention} & Vaswani et al. (2017) & Multi-Head Self-Attention Seq2Seq Transformer & WMT14 En-De (4.5M) and En-Fr (36M pairs) & Eliminates recurrence; enables massive GPU parallelization; models long-range syntax effectively & \textbf{Drawback:} Quadratic memory $O(N^2)$; high data hunger; severely overfits on low-resource Kannada without massive pretraining. \textit{Addressed by using Indic-pretrained models with Dravidian subword tokenization.} \\
\addlinespace[0.5ex]

\cite{haddow2022survey} & Haddow et al. (2022) & Comprehensive Survey of Low-Resource MT & Global low-resource and underrepresented pairs & Establishes taxonomy of bitext mining, back-translation, and multilingual transfer & \textbf{Drawback:} Methodological survey only; lacks production architecture or Kannada-specific empirical pipeline. \textit{Addressed by realizing an end-to-end verified news ingestion and translation system.} \\
\addlinespace[0.5ex]

\cite{liu2020multilingual} & Liu et al. (2020) & Multilingual Denoising Seq2Seq (mBART) & CC25 corpus (25 languages, Common Crawl) & Pretrain-then-finetune paradigm yields large gains on resource-constrained pairs & \textbf{Drawback:} \textbf{Kannada completely omitted} from pretraining; unable to process Kannada script without vocabulary surgery. \textit{Addressed by adopting native Indic-specific architectures.} \\
\addlinespace[0.5ex]

\cite{tang2020multilingual} & Tang et al. (2020) & Incremental Continual Pretraining (mBART-50) & ML50 multilingual benchmark (50 languages) & Direct many-to-many translation without English pivot; robust multilingual transfer & \textbf{Drawback:} \textbf{Kannada omitted} from the 50 languages; tokenizer lacks Kannada subword representations. \textit{Addressed by employing models trained explicitly on the 22 scheduled Indian languages.} \\
\addlinespace[0.5ex]

\cite{nllb2024scaling} & NLLB Team (2024) & Sparsely-Gated Mixture-of-Experts (NLLB-200) & 200 languages; mined web corpora; FLORES-200 & Massive language coverage; native Kannada support; high benchmark BLEU/chrF++ & \textbf{Drawback:} Heavy parameter size ($>$3B); high inference latency; general web data causes named-entity drift and hallucinations in regional news. \textit{Addressed by lightweight sentence-aligned failovers with quality guards.} \\
\addlinespace[0.5ex]

\cite{ramesh2022samanantar} & Ramesh et al. (2022) & Automated Parallel Bitext Mining (Samanantar) & 11 Indic languages + English (49.7M pairs) & Largest open parallel Indic dataset; 4M+ clean Kannada-English sentence pairs & \textbf{Drawback:} General web-crawled noise; high domain gap when applied to localized regional journalism and headlines. \textit{Addressed by domain-tailored sentence normalization and news-specific quality estimation.} \\
\addlinespace[0.5ex]

\cite{gala2023indictrans2} & Gala et al. (2023) & Indic Pretrained Encoder-Decoder (IndicTrans2) & All 22 scheduled Indian languages (BPCC, IN22) & SOTA translation on Indic languages; script unification; native Dravidian tokenization & \textbf{Drawback:} Pure translation model; lacks automated news ingestion, quality guardrails, summarization, and verifiable citations. \textit{Addressed by building an end-to-end verified Translate-Then-Summarize pipeline around it.} \\
\addlinespace[0.5ex]

\cite{goyal2022flores} & Goyal et al. (2022) & Multi-Way Parallel Evaluation (FLORES-101) & 101 languages, 3010 translated sentences & High-quality professional translations; strict multi-way comparability & \textbf{Drawback:} Limited to Wikipedia domain; completely lacks journalistic syntax, dynamic political jargon, and breaking news style. \textit{Addressed by establishing evaluation on live Kannada newspaper feeds.} \\
\addlinespace[0.5ex]

\cite{rei2020comet} & Rei et al. (2020) & Cross-Lingual Metric Predictor (COMET) & WMT Metrics Shared Task human datasets & Superior correlation with human judgments; evaluates semantic adequacy beyond n-grams & \textbf{Drawback:} Requires reference translations; high computational overhead; cannot serve as a live inference guardrail. \textit{Addressed by implementing lightweight reference-free Quality Estimation (QE) metrics.} \\
\addlinespace[0.5ex]

\cite{hendy2023gpt} & Hendy et al. (2023) & Prompt-Based In-Context Learning (GPT-3.5/GPT-4) & 18 translation directions across diverse tiers & Competitive on high-resource pairs; strong fluency, post-editing, and abstractive synthesis & \textbf{Drawback:} Frequent hallucinations on low-resource Indic languages; drops named entities; lacks sentence traceability. \textit{Addressed by strictly separating sentence translation from closed-book attributed summarization.} \\
\bottomrule
\end{tabularx}
\end{table*}

% -----------------------------------------------------------------------------
% SECTION III: PROPOSED METHODOLOGY AND SYSTEM ARCHITECTURE
% -----------------------------------------------------------------------------
\section{Proposed Methodology and System Architecture}
\label{sec:methodology}

The proposed system addresses the limitations of unconstrained generative models by implementing a multi-stage, modular \textit{Translate-Then-Summarize} pipeline fortified with multi-tiered quality guardrails.

\subsection{Overall System Architecture}
The system architecture comprises five decoupled sub-systems:
\begin{enumerate}
    \item \textbf{Multi-Modal Ingestion Engine:} Ingests Kannada news from live RSS feeds, XML News Sitemaps, digital/scanned PDFs, and e-paper image clippings.
    \item \textbf{Linguistic Preprocessing \& Segmentation:} Performs Unicode canonical normalization, boilerplate noise filtration, and Sentence Boundary Detection (SBD).
    \item \textbf{Sentence-Aligned Translation Engine:} Executes sentence-by-sentence neural translation with $1$-to-$1$ bidirectional mapping and graceful failover.
    \item \textbf{Quality Estimation (QE) \& Hallucination Guardrails:} Audits sentence translations using token length bounds, numeric invariant tracking, and cross-lingual semantic similarity.
    \item \textbf{Grounded Summarization \& Verifiable Presentation:} Produces evidence-attributed summaries, powers an interactive dual-reader, enables sentence-synchronized TTS briefings, and facilitates grounded Q\&A.
\end{enumerate}

\subsection{Multi-Modal Ingestion Engine}
To accommodate the diverse formats in which regional news is published across Karnataka, the ingestion layer provides three specialized pipelines:

\subsubsection{Automated RSS and XML News Sitemap Aggregator}
The aggregator connects to eight premier Kannada news outlets, categorized into standard RSS feeds and XML Google News sitemaps:
\begin{itemize}
    \item \textbf{Direct RSS Feeds:} \textit{Prajavani}, \textit{Vijay Karnataka}, \textit{News18 Kannada}, \textit{Asianet News Kannada}, and \textit{Times Now Kannada}.
    \item \textbf{XML News Sitemaps:} \textit{Kannada Prabha} (via dynamic news sitemaps containing Google News XML tags) and \textit{Udayavani} (via monthly hierarchical XML indexes).
    \item \textbf{Single-Page Application (SPA) Dynamic Asset Extractor:} Certain modern portals (e.g., \textit{Udayavani}) render content dynamically on the client side via Content Delivery Networks (CDNs). The aggregator detects CDN JSON endpoints embedded in Next.js payloads (e.g., \texttt{/test-news/news-json/*.json}), fetching raw story objects and stripping styling markup to recover the clean source text.
\end{itemize}
All outgoing requests enforce compliance with \texttt{robots.txt} policies and integrate randomized exponential backoff to ensure ethical web scraping.

\subsubsection{Digital and Scanned PDF Ingestion Engine}
For official government orders, state gazettes, and press releases distributed as PDF documents, the ingestion module utilizes a dual-path extractor:
\begin{itemize}
    \item \textbf{Digital Extraction:} Native text streams are parsed page-by-page using \texttt{pypdf}. Whitespace, running headers, and pagination footers are algorithmically removed.
    \item \textbf{Image Fallback Extraction:} If the digital character yield is below a predefined threshold ($\text{len}(text) < 60$), the document is identified as a scanned facsimile. The system extracts embedded page raster images and routes them to the Vision OCR pipeline.
\end{itemize}

\subsubsection{E-Paper and News Screenshot Vision OCR Engine}
Physical newspaper clippings, cropped digital e-paper columns, and smartphone news screenshots are processed using specialized multimodal vision encoders. The raw image is normalized, base64-encoded, and ingested by Gemini Vision models with a strict domain-specific system prompt directing the model to transcribe Kannada Unicode text verbatim while preserving paragraph hierarchy and eliminating extraneous advertisement text.

\subsection{Linguistic Preprocessing and Sentence Boundary Detection}
Raw Kannada text extracted from diverse sources contains typographical inconsistencies, non-standard diacritics, and irregular punctuation. The preprocessing layer executes:
\begin{enumerate}
    \item \textbf{Unicode Normalization:} Kannada script representations are normalized to Unicode Canonical Decomposition followed by Canonical Composition (NFC), resolving inconsistent encodings of Kannada vowel signs (e.g., \textit{ottu} forms and dependent vowel signs like \textit{dirgha}).
    \item \textbf{Boilerplate and Ad Filtration:} A regex-based domain filter purges ubiquitous news artifacts, including subscription appeals (``\textit{Chandaa daaranagi}''), social media sharing prompts (``\textit{Follow maadi}''), advertisement placeholders (``\textit{Jaahiraatu}''), and reporter bylines.
    \item \textbf{Predominant Script Verification:} To prevent processing corrupt inputs, a script ratio metric calculates:
    \begin{equation}
        \mathcal{R}_{kn} = \frac{\sum_{c \in T} \mathbb{I}(0x0C80 \le \text{ord}(c) \le 0x0CFF)}{\sum_{c \in T} \mathbb{I}(c \in \text{Alphabetic})}
        \label{eq:script_ratio}
    \end{equation}
    where $T$ represents the input character sequence. Articles with $\mathcal{R}_{kn} < 0.40$ are rejected.
    \item \textbf{Sentence Boundary Detection (SBD):} Kannada journalistic text employs both Western punctuation (periods, question marks) and traditional Indic punctuation (\textit{danda} U+0964, \textit{double danda} U+0965). The SBD algorithm splits text across these boundaries while guarding against false splits on abbreviated honorifics (e.g., \textit{Dr.}, \textit{Shri.}) and numeric decimals:
    \begin{equation}
        S_{kn} = \{s_0^{kn}, s_1^{kn}, s_2^{kn}, \dots, s_{N-1}^{kn}\}
        \label{eq:sentence_set}
    \end{equation}
\end{enumerate}

\subsection{Sentence-Level Neural Translation and Alignment}
Rather than submitting an entire news article as an unconstrained monolithic block, the translation engine operates strictly over the ordered sentence sequence $S_{kn}$. This guarantees fine-grained sentence alignment and prevents the omission or hallucination of complete paragraphs.

Let $f_{\theta}: \mathcal{X}_{kn} \to \mathcal{Y}_{en}$ denote the neural translation mapping parameterized by weights $\theta$. For each source sentence $s_i^{kn} \in S_{kn}$, the engine generates:
\begin{equation}
    s_i^{en} = f_{\theta}(s_i^{kn}) \quad \forall i \in \{0, 1, \dots, N-1\}
    \label{eq:sentence_translation}
\end{equation}
The primary model utilized is \textbf{IndicTrans2} (1B parameter Indic-to-English variant), which incorporates specialized tokenizers tailored to Dravidian morphophonemic structures. 

To guarantee high availability and prevent rate-limiting during bulk ingestion, the system implements an automated \textbf{Multi-Tier Failover Architecture}:
\begin{itemize}
    \item \textbf{Tier 1 (Specialized Indic MT):} IndicTrans2 executed locally or via dedicated inference endpoints.
    \item \textbf{Tier 2 (Batch Generative MT):} High-speed multimodal LLMs (\texttt{gemini-3.1-flash-lite} / \texttt{gemini-3-flash-preview}) queried with a strict array-preserving batch translation prompt.
    \item \textbf{Tier 3 (Web NMT Fallback):} Google Translate web endpoint executed sentence-by-sentence with exponential backoff.
\end{itemize}
Regardless of the active tier, the output preserves an invariant $1$-to-$1$ bijection between $s_i^{kn}$ and $s_i^{en}$.

\subsection{Quality Estimation and Hallucination Guardrails}
Each translated sentence pair $(s_i^{kn}, s_i^{en})$ is immediately subjected to four reference-free Quality Estimation (QE) guardrails:

\subsubsection{Token Length Ratio Guard}
Because Kannada is morphologically agglutinative and English is analytic, the ratio of English target tokens to Kannada source tokens adheres to a well-defined statistical distribution:
\begin{equation}
    \gamma_{len}(s_i) = \frac{|\text{Tokens}_{en}(s_i^{en})|}{|\text{Tokens}_{kn}(s_i^{kn})|}
    \label{eq:length_ratio}
\end{equation}
Under normal translation, $1.1 \le \gamma_{len} \le 2.6$. If $\gamma_{len} < 0.60$ (indicating severe under-translation or sentence truncation) or $\gamma_{len} > 3.80$ (indicating runaway degenerative repetition or hallucination), a flag $\mathcal{F}_{length}$ is raised.

\subsubsection{Numeric and Date Invariant Preserver}
News journalism depends fundamentally on numeric precision (casualties, financial figures, dates, legislative bill numbers). Let $\mathcal{N}(s)$ denote the set of numeric entities extracted from sentence $s$:
\begin{equation}
    \mathcal{N}_{kn}(s_i^{kn}) \longleftrightarrow \mathcal{N}_{en}(s_i^{en})
    \label{eq:numeric_preservation}
\end{equation}
The guard converts all Kannada numerals (U+0CE6 to U+0CEF: ೦, ೧, ೨, ೩, ೪, ೫, ೬, ೭, ೮, ೯) to standard ASCII Arabic digits and verifies set equivalence. If $\mathcal{N}_{kn} \setminus \mathcal{N}_{en} \ne \emptyset$, a numeric distortion flag $\mathcal{F}_{num}$ is recorded.

\subsubsection{Named-Entity and Transliteration Audit}
High-frequency regional named entities $\mathcal{E}$ (e.g., proper nouns, political figures, geographic locations) are extracted using a rule-assisted Named Entity Recognizer. The guard confirms that for every entity $e \in \mathcal{E}_{kn}$, its corresponding English phoneme/transliteration appears in $s_i^{en}$.

\subsubsection{Cross-Lingual Semantic Similarity via LaBSE}
To verify semantic adequacy without reference translations, both sentences are projected into a common 768-dimensional multilingual vector space using Language-Agnostic BERT Sentence Embeddings (LaBSE)~\cite{feng2022language}:
\begin{equation}
    \mathbf{v}_{kn} = \text{LaBSE}(s_i^{kn}), \quad \mathbf{v}_{en} = \text{LaBSE}(s_i^{en})
    \label{eq:labse_embeddings}
\end{equation}
The semantic alignment score is computed via cosine similarity:
\begin{equation}
    \mathcal{S}_{sim}(s_i^{kn}, s_i^{en}) = \frac{\mathbf{v}_{kn} \cdot \mathbf{v}_{en}}{\|\mathbf{v}_{kn}\|_2 \|\mathbf{v}_{en}\|_2}
    \label{eq:cosine_sim}
\end{equation}
Sentence translations exhibiting $\mathcal{S}_{sim} < 0.65$ are marked as suspicious ($\mathcal{F}_{sim}$), automatically triggering an article status change from \texttt{ready} to \texttt{needs\_review}.

\subsection{Grounded Abstractive Summarization with Attribution}
Following verified sentence-level translation, the system synthesizes an executive journalistic summary. To eliminate hallucination risks, the summarization engine is structured as a \textbf{Closed-Book Grounded Synthesizer}.

The prompt context is strictly bounded to the numbered English sentence sequence:
\begin{equation}
    \mathcal{C} = \big\{ [S_0]\; s_0^{en}, \; [S_1]\; s_1^{en}, \; \dots, \; [S_{N-1}]\; s_{N-1}^{en} \big\}
    \label{eq:context_prompt}
\end{equation}
The summarizer is constrained via JSON schema enforcement to output:
\begin{itemize}
    \item \textbf{Headline ($\mathcal{H}_{en}$):} A concise, objective English headline.
    \item \textbf{Two-Line Executive Summary ($\Sigma_{brief}$):} A two-sentence synthesis of the core development.
    \item \textbf{Key Fact Points ($\mathcal{K}$):} A list of 3--5 bullet statements $\{k_1, k_2, \dots, k_M\}$. Crucially, each bullet point must explicitly cite its supporting evidence sentence indices:
    \begin{equation}
        k_j = \big( \text{Text}_j, \; \mathcal{E}_j \subseteq \{0, 1, \dots, N-1\} \big)
        \label{eq:key_point_evidence}
    \end{equation}
\end{itemize}

\subsubsection{7-Layer Faithfulness Verification and Extractive Fallback}
Before committing the summary to the database, an automated verification engine evaluates the generated key points:
\begin{enumerate}
    \item \textbf{Index Range Validation:} Confirms that all cited indices $\forall e \in \mathcal{E}_j$ satisfy $0 \le e < N$.
    \item \textbf{Lexical Entailment Scoring:} Evaluates whether the tokens comprising $\text{Text}_j$ demonstrate substantive n-gram overlap and semantic entailment with the cited evidence sentences $\bigcup_{e \in \mathcal{E}_j} s_e^{en}$.
    \item \textbf{Extractive Fallback:} If any key point cites invalid indices or exhibits semantic divergence exceeding tolerance thresholds, the generative summary is discarded. In its place, the system activates an algorithmic TextRank extractive summarizer that extracts the top 3 central sentences directly from $S_{en}$, ensuring absolute zero-hallucination guarantees.
\end{enumerate}

\subsection{Interactive Multimodal Presentation \& Conversational Q\&A}
The translated and verified output is rendered through a Next.js web application designed for accessibility, transparency, and deep inquiry:
\begin{itemize}
    \item \textbf{Interactive Dual-Reader:} Kannada source paragraphs and English translations are displayed in synchronized side-by-side columns. Hovering over or selecting any sentence in either column dynamically highlights its corresponding translation in the opposing language.
    \item \textbf{Evidence Pill Traversal:} Clicking an evidence badge (e.g., \texttt{[S2]}) on any summary key point instantly scrolls the dual-reader to sentence 2 and highlights both the Kannada and English source text.
    \item \textbf{Speed-Adaptive Audio Briefing:} An embedded Web Speech synthesis engine delivers an English audio briefing of the headline, executive summary, and key points. The player utilizes an active sentence queue architecture: modifying the speech playback rate ($1\times, 1.25\times, 1.5\times, 2\times$) dynamically updates the rate of the currently active sentence without resetting playback to the beginning.
    \item \textbf{``Ask the Article'' Grounded Chatbot:} Users can submit arbitrary queries regarding the news story. The query is resolved against the numbered sentence context $\mathcal{C}$ using an extractive fallback and strict prompt boundaries, returning concise answers paired with clickable evidence citations.
\end{itemize}

% -----------------------------------------------------------------------------
% REFERENCES
% -----------------------------------------------------------------------------
\begin{thebibliography}{00}

\bibitem{vaswani2017attention}
A.~Vaswani, N.~Shazeer, N.~Parmar, J.~Uszkoreit, L.~Jones, A.~N. Gomez, {\L}.~Kaiser, and I.~Polosukhin, ``Attention is all you need,'' in \textit{Advances in Neural Information Processing Systems (NeurIPS)}, vol.~30, 2017, pp. 5998--6008.

\bibitem{haddow2022survey}
B.~Haddow, R.~Bawden, A.~V. Miceli~Barone, J.~Helcl, and A.~Birch, ``Survey of low-resource machine translation,'' \textit{Computational Linguistics}, vol.~48, no.~3, pp. 673--732, Sep. 2022.

\bibitem{liu2020multilingual}
Y.~Liu, J.~Gu, N.~Goyal, X.~Li, S.~Edunov, M.~Ghazvininejad, M.~Lewis, and L.~Zettlemoyer, ``Multilingual denoising pre-training for neural machine translation,'' \textit{Transactions of the Association for Computational Linguistics (TACL)}, vol.~8, pp. 726--742, 2020.

\bibitem{tang2020multilingual}
Y.~Tang, C.~Tran, X.~Li, P.-J. Chen, N.~Goyal, V.~Chaudhary, J.~Gu, and A.~Fan, ``Multilingual translation with extensible multilingual pretraining and finetuning,'' \textit{arXiv preprint arXiv:2008.00401}, 2020.

\bibitem{nllb2024scaling}
{NLLB Team}, M.~R. Costa-juss{\`a}, J.~Cross, O.~Celebi, M.~Elbayad, K.~Heafield, K.~Heffernan, E.~Kalyanpur, J.~Schwenk, et al., ``Scaling neural machine translation to 200 languages,'' \textit{Nature}, vol.~630, pp. 841--846, 2024.

\bibitem{ramesh2022samanantar}
G.~Ramesh, S.~Doddapaneni, A.~Bapna, M.~M. Khapra, and P.~Kumar, ``Samanantar: The largest publicly available parallel corpora collection for 11 Indic languages,'' \textit{Transactions of the Association for Computational Linguistics (TACL)}, vol.~10, pp. 145--162, 2022.

\bibitem{gala2023indictrans2}
J.~Gala, P.~A. Chitale, R.~K. Raghavan, V.~Joshi, S.~Doddapaneni, V.~Rangan, A.~M. P., A.~K. B., et al., ``IndicTrans2: Towards high-quality and accessible machine translation models for all 22 scheduled Indian languages,'' \textit{Transactions on Machine Learning Research (TMLR)}, 2023.

\bibitem{goyal2022flores}
N.~Goyal, S.~Gao, V.~Chaudhary, P.-J. Chen, G.~Wenzek, D.~Ju, S.~Krishnan, M.~Shen, N.~Lakhotia, et al., ``The Flores-101 evaluation benchmark for low-resource and multilingual machine translation,'' \textit{Transactions of the Association for Computational Linguistics (TACL)}, vol.~10, pp. 522--538, 2022.

\bibitem{rei2020comet}
R.~Rei, C.~Stewart, A.~C. Farinha, and A.~Lavie, ``COMET: A neural framework for MT evaluation,'' in \textit{Proc. Conf. Empirical Methods in Natural Language Processing (EMNLP)}, 2020, pp. 2685--2702.

\bibitem{hendy2023gpt}
A.~Hendy, M.~Abdelrehim, A.~Amr, M.~Abdelghany, M.~Khaled, M.~Kamal, M.~Hany, R.~Mohey, et al., ``How good are GPT models at machine translation? A comprehensive evaluation,'' \textit{arXiv preprint arXiv:2302.09210}, 2023.

\bibitem{feng2022language}
F.~Feng, Y.~Yang, D.~Cer, N.~Arivazhagan, and W.~Wang, ``Language-agnostic BERT sentence embeddings,'' in \textit{Proc. 60th Annual Meeting of the Association for Computational Linguistics (ACL)}, 2022, pp. 878--891.

\end{thebibliography}

\end{document}
